\begin{abstract}
When a language model continues a reasoning trace that already contains a
false step, does it catch the error or quietly build on it? We study this
with an audited perturb-and-validate instrument: a statement is planted into
a model's own proof trace on formally checkable logic problems, audited for
whether it is true or false under the closed world, and the model's
unprompted continuation is formally validated.

Responses to a planted false step are strongly model-dependent. On the
primary model, Qwen2.5-7B, explicit rejection is a visibility event: a
planted falsehood is rejected in roughly 0.25 of continuations when its
contradiction is a directly readable premise and in about 0.03 one hop away,
flat beyond, and genuine multi-hop checking exists but is rare and ungraded
by proximity. This visible-contradiction rejection pattern does not transfer:
on three of four other open-weights models it significantly reverses, with
rejection \emph{lower} when the contradiction is visible than when it is a
hop away.

Among the models that absorb planted falsehoods, evidence distance does not
predict absorption but derivational usefulness does: a usable falsehood is
reused near 0.8 regardless of how close its refutation sits, an inert one
near zero. Absorption itself declines to near zero with scale within one
family (0.81 at 7B to 0.01 at 32B in Qwen2.5) and varies by training recipe
across labs at a matched harness.

Two results hold across scale and lineage. No tested model re-verifies an
already-computed value: recomputable falsehoods are absorbed near 1.00 even
by frontier models that reject re-readable ones. And verbal rejection is a
weak proxy for checking: a frequency-matched non-refuting rule elicits the
same rejection rate as genuine refutation, yet those rejections fabricate a
rule absent from the world, so only a derivation-licensed split separates
checking from lexical cueing.

Single-model conclusions about how reasoning errors are handled do not
transfer across model, scale, or training recipe; computed intermediate
values are the one universal blind spot. We frame the in-stream-versus-fresh-turn
part of this last boundary as corroboration of concurrent work on the
self-correction blind spot, not as a discovery, and our controlled tests run
on a single model family and a synthetic logic domain.
\end{abstract}
